\documentclass[lettersize,journal]{IEEEtran}
\usepackage{amsmath,amsfonts}
\usepackage{algorithmic}
\usepackage{algorithm}
\usepackage{array}
\usepackage[caption=false,font=normalsize,labelfont=sf,textfont=sf]{subfig}
\usepackage{textcomp}
\usepackage{stfloats}
\usepackage{url}
\usepackage{verbatim}
\usepackage{graphicx}
\usepackage{cite}
\hyphenation{op-tical net-works semi-conduc-tor IEEE-Xplore}

\begin{document}

\title{Delay-Robust Closed-Loop Inverse Kinematics Control for Unified 25-DoF Arm-Hand Dexterous Teleoperation Systems}

\author{Thien-Bao Tran, and Bao-Tran Nguyen
\thanks{Manuscript received July 21, 2026; revised July 21, 2026. This work was supported by Advanced Robotic Manipulation Lab.}}

\markboth{IEEE Transactions on Robotics,~Vol.~22, No.~7, July~2026}%
{Tran \MakeLowercase{\textit{et al.}}: Delay-Robust Closed-Loop Inverse Kinematics Control for Unified 25-DoF Arm-Hand Systems}

\maketitle

\begin{abstract}
This paper presents a delay-robust Closed-Loop Inverse Kinematics (CLIK) control scheme for a unified 25-DoF arm-hand robotic system (6-DoF UR5 manipulator and 19-DoF custom dexterous hand) under time-varying communication delays. By introducing a filtered reference trajectory and modeling latency under a true-and-noise framework, we prove Local Input-to-State Stability (LISS) using a Lyapunov-Krasovskii Functional (LKF) and Linear Matrix Inequalities (LMIs). A damped least-squares (DLS) Jacobian inverse is utilized to guarantee singularity-robust tracking, while secondary null-space projection optimizes joint limits, self-collision avoidance, and natural anthropomorphic postures. The effectiveness of the proposed method is validated at 100 Hz in a physical simulator.
\end{abstract}

\begin{IEEEkeywords}
Closed-Loop Inverse Kinematics, Lyapunov-Krasovskii Functional, Input-to-State Stability, Teleoperation, Redundant Robots.
\end{IEEEkeywords}

\section{Introduction}
\IEEEPARstart{V}{ision-based} teleoperation systems enable human operators to control high-DoF robotic manipulators in real-time. However, time-varying communication latencies and packet jitter often compromise the closed-loop tracking performance. This paper proposes a unified CLIK framework to guarantee stability and collision avoidance under variable delay.

\section{Related Work}
Existing time-delay teleoperation algorithms often neglect the kinematic singularities and highly redundant joint configurations of integrated arm-hand systems. We bridge this gap by formulating a 25-DoF unified Jacobian with LMI-based stability constraints.

\section{Problem Formulation}

Consider an integrated robotic system consisting of a 6-DoF manipulator (UR5) and a 19-DoF custom dexterous hand, constituting a highly redundant kinematic chain with a joint configuration vector $q(t) \in \mathbb{R}^n$ ($n = 25$), defined as:
\begin{equation}
q(t) = \begin{bmatrix} q_a(t) \\ q_h(t) \end{bmatrix}
\end{equation}
where $q_a(t) \in \mathbb{R}^6$ represents the joint angles of the robotic arm, and $q_h(t) \in \mathbb{R}^{19}$ denotes the active joint angles of the multi-fingered hand. 

The forward kinematics mapping the joint space configuration to the task space pose $x(t) \in \mathcal{X}$ is given by:
\begin{equation}
x(t) = f(q(t))
\end{equation}
where the task space $\mathcal{X} = \mathbb{R}^3 \times \mathcal{S}^3 \times \mathbb{R}^{n_f}$ is designed with $n_f = 5$ relative Euclidean finger opposition distances (specifically, thumb-index, thumb-middle, thumb-ring, thumb-pinky, and index-pinky opposition pairs) to prevent kinematic over-constraint ($m = 3 + 3 + 5 = 11$ task-space dimensions, which is well below $n = 25$ total joints). 

Differentiating the forward kinematics yields the differential kinematics:
\begin{equation}
v(t) = J(q(t))\dot{q}(t)
\end{equation}
where $v(t) = [\dot{p}_w^T(t), \omega_w^T(t), \dot{d}_f^T(t)]^T \in \mathbb{R}^m$ is the spatial task-space velocity vector ($\omega_w$ being the physical angular velocity of the wrist), and $J(q(t)) \in \mathbb{R}^{m \times n}$ is the analytical system Jacobian.

In a remote teleoperation setup, the desired trajectory $x_d(t)$ is commanded by a human operator through a vision-based hand tracker. The controller receives $x_d(t)$ under an unknown but deterministic time-varying communication delay $\tau(t)$, satisfying the following conditions:

\textbf{Assumption 1}: The communication delay $\tau(t)$ is continuous and bounded such that:
\begin{equation}
0 \le \tau(t) \le \tau_m, \quad \forall t \ge 0
\end{equation}
where $\tau_m$ is a known upper bound of the latency.

\textbf{Assumption 2}: The rate of change of the true delay is bounded on both sides to prevent unstable tracking:
\begin{equation}
|\dot{\tau}_{true}(t)| \le d < 1, \quad \forall t \ge 0
\end{equation}

\textbf{Assumption 3}: The measured latency consists of a true delay $\tau_{true}(t)$ satisfying Assumptions 1 and 2, and a high-frequency bounded measurement noise $n(t)$ representing non-differentiable network packet jitter:
\begin{equation}
\tau(t) = \tau_{true}(t) + n(t), \quad \|n(t)\| \le n_{max}
\end{equation}

Because the raw delayed trajectory $x_d(t - \tau(t))$ contains the high-frequency non-differentiable jitter $n(t)$, directly feeding it into a continuous-time differential controller violates smoothness requirements. In practice, the reference signal is processed through a continuous-time low-pass filter (or jitter buffer). Thus, we define the \textit{filtered reference trajectory} $x_{d,filt}(t)$ equivalent to evaluating the trajectory at the filtered delay $\tau_{filt}(t)$:
\begin{equation}
x_{d,filt}(t) = x_d(t - \tau_{filt}(t))
\end{equation}
where $\tau_{filt}(t)$ is generated such that it is smooth, strictly bounded ($0 \le \tau_{filt}(t) \le \tau_m$), and its derivative is strictly bounded ($|\dot{\tau}_{filt}(t)| \le d < 1$).

Under delayed teleoperation, the task-space tracking error $e(t) \in \mathbb{R}^m$ is redefined relative to the filtered reference:
\begin{equation}
e(t) = \begin{bmatrix} e_p(t) \\ e_o(t) \\ e_f(t) \end{bmatrix} = \begin{bmatrix} p_{d}(t-\tau_{filt}(t)) - p_w(t) \\ \eta_d(t-\tau_{filt}(t))\epsilon(t) - \eta(t)\epsilon_d(t-\tau_{filt}(t)) - \epsilon^\times(t)\epsilon_d(t-\tau_{filt}(t)) \\ d_{d}(t-\tau_{filt}(t)) - d_f(t) \end{bmatrix}
\end{equation}
where $e_o(t) \in \mathbb{R}^3$ represents the quaternion-based orientation error, and $\epsilon^\times$ is the skew-symmetric matrix of $\epsilon$. By redefining the tracking error relative to the filtered reference, the physical jitter $n(t)$ is decoupled from the closed-loop differential kinematics, acting instead as a bounded spatial tracking offset (since $x_d$ is Lipschitz continuous and $n(t)$ is bounded by Assumption 3).

The primary control objective is to design a joint velocity controller $\dot{q}(t)$ that guarantees Input-to-State Stability (ISS) of the tracking error under Assumptions 1, 2, and 3, while utilizing the null-space to safely avoid joint limits, self-collisions, and unnatural hand configurations.

\section{Proposed Scheme}

\subsection{Singularity-Robust CLIK Control Law}
To prevent numerical instability near kinematic singularities, we employ a Damped Least-Squares (DLS) Jacobian inverse. The proposed control law is formulated as:
\begin{equation}
\dot{q}(t) = J^{\dagger}_{DLS}(q(t)) \left[ v_{d,filt}(t) + K_p e(t) + K_d e(t - \tau_{filt}(t)) \right] + \left( I_n - J^{\dagger}_{DLS}(q(t))J(q(t)) \right) \dot{q}_0(t)
\end{equation}
where:
\begin{itemize}
\item $J^{\dagger}_{DLS}(q) = J^T(JJ^T + \lambda^2 I_m)^{-1} \in \mathbb{R}^{n \times m}$ is the DLS pseudo-inverse, with $\lambda > 0$ being the damping factor.
\item $K_p, K_d \in \mathbb{R}^{m \times m}$ are positive-definite diagonal gain matrices.
\item $K_d e(t - \tau_{filt}(t))$ is the discrete delay-feedback compensation term.
\item $v_{d,filt}(t) = \frac{d}{dt} x_d(t - \tau_{filt}(t)) = v_d(t - \tau_{filt}(t))(1 - \dot{\tau}_{filt}(t))$ is the analytical derivative of the filtered reference trajectory.
\item $\dot{q}_0(t) \in \mathbb{R}^n$ is the joint velocity vector designed for secondary task optimization.
\end{itemize}

Substituting the control law into the differential kinematics yields the closed-loop tracking error dynamics:
\begin{equation}
\dot{e}(t) = -E(e_o) \left[ K_p e(t) + K_d e(t - \tau_{filt}(t)) \right] + d_t(t)
\end{equation}
where $E(e_o) = \text{diag}\left(I_3, \frac{1}{2}(\eta_e I_3 + \epsilon_e^\times), I_{5}\right)$ is the orientation scaling matrix, and $d_t(t)$ is the aggregate perturbation term:
\begin{equation}
d_t(t) = d_s(t) + d_l(t)
\end{equation}
with:
\begin{itemize}
\item $d_s(t) = \left( J J^{\dagger}_{DLS} - I_m \right) \left[ v_{d,filt}(t) + K_p e(t) + K_d e(t - \tau_{filt}(t)) \right]$ being the damping perturbation.
\item $d_l(t) = -J(q)\left( I_n - J^{\dagger}_{DLS}(q)J(q) \right)\dot{q}_0(t)$ being the workspace leakage originating from the DLS null-space projection.
\end{itemize}

To rigorously address the nonlinearities in $E(e_o)$, we define the linearization residual matrix mapping $\Delta E(e_o) = I_m - E(e_o)$. This allows the exact closed-loop error kinematics to be formulated as:
\begin{equation}
\dot{e}(t) = -K_p e(t) - K_d e(t - \tau_{filt}(t)) + d_t(t) + \Delta E(e_o) \left[ K_p e(t) + K_d e(t - \tau_{filt}(t)) \right]
\end{equation}

\textbf{Remark 1 (Local ISS)}: Within the region of attraction $\Omega_o = \{ e_o \in \mathbb{R}^3 \mid \eta_e > 0 \}$ (corresponding to rotation errors of less than $180^\circ$), the residual mapping $\Delta E(e_o) = I_m - E(e_o)$ converges to zero near the equilibrium state. By restricting the stability analysis to a local neighborhood where $\|\Delta E(e_o)\|$ is sufficiently small, its higher-order perturbative effects on the LKF derivative are strictly dominated by the negative-definite matrix $\Omega_{ISS}$. Thus, the established stability property is Local Input-to-State Stability (Local ISS).

\subsection{Stability Analysis via Lyapunov-Krasovskii Functional}
To prove the Local ISS of the closed-loop error system, we construct the following candidate Lyapunov-Krasovskii Functional (LKF) $V(t)$:
\begin{equation}
V(t) = e^T(t) P e(t) + \int_{t - \tau_{filt}(t)}^{t} e^T(s) Q e(s) ds + \tau_m \int_{-\tau_m}^{0} \int_{t + \theta}^{t} \dot{e}^T(s) R \dot{e}(s) ds d\theta
\end{equation}
where $P, Q, R \in \mathbb{R}^{m \times m}$ are diagonal positive-definite matrices (hence $P K_p = K_p P$, $P K_d = K_d P$).

Differentiating $V(t)$ along the trajectory of the perturbed error dynamics yields:
\begin{equation}
\dot{V}(t) = 2 e^T(t) P \left[ -K_p e(t) - K_d e(t - \tau_{filt}(t)) + d_t(t) \right] + e^T(t) Q e(t) - (1 - \dot{\tau}_{filt}(t)) e^T(t - \tau_{filt}(t)) Q e(t - \tau_{filt}(t)) + \tau_m^2 \dot{e}^T(t) R \dot{e}(t) - \tau_m \int_{t - \tau_m}^{t} \dot{e}^T(s) R \dot{e}(s) ds
\end{equation}

Using Young's inequality, the bilinear perturbation term is bounded by:
\begin{equation}
2 e^T(t) P d_t(t) \le \epsilon_1 e^T(t) e(t) + \frac{1}{\epsilon_1} d_t^T(t) P^2 d_t(t)
\end{equation}
where $\epsilon_1 > 0$.

Let $\phi_e(t) = -K_p e(t) - K_d e(t - \tau_{filt}(t))$. The acceleration-related LKF term contains:
\begin{equation}
\dot{e}^T(t) R \dot{e}(t) = \left[ \phi_e(t) + d_t(t) \right]^T R \left[ \phi_e(t) + d_t(t) \right]
\end{equation}
Applying Young's inequality to the cross-term $2 \phi_e^T(t) R d_t(t)$ yields:
\begin{equation}
\dot{e}^T(t) R \dot{e}(t) \le (1 + \epsilon_2) \phi_e^T(t) R \phi_e(t) + \left( 1 + \frac{1}{\epsilon_2} \right) d_t^T(t) R d_t(t)
\end{equation}
where $\epsilon_2 > 0$.

To handle the delay-dependent integral term, we partition the integration interval into two segments: $[t - \tau_m, t - \tau_{filt}(t)]$ and $[t - \tau_{filt}(t), t]$. Applying Jensen's inequality to each segment yields:
\begin{equation}
-\tau_m \int_{t - \tau_{filt}(t)}^{t} \dot{e}^T(s) R \dot{e}(s) ds \le -\frac{\tau_m}{\tau_{filt}(t)} a^T(t) R a(t)
\end{equation}
\begin{equation}
-\tau_m \int_{t - \tau_m}^{t-\tau_{filt}(t)} \dot{e}^T(s) R \dot{e}(s) ds \le -\frac{\tau_m}{\tau_m - \tau_{filt}(t)} b^T(t) R b(t)
\end{equation}
where $a(t) = e(t) - e(t-\tau_{filt}(t))$ and $b(t) = e(t-\tau_{filt}(t)) - e(t-\tau_m)$.

Let $\eta(t) = [e^T(t), e^T(t-\tau_{filt}(t)), e^T(t-\tau_m)]^T \in \mathbb{R}^{3m}$ be the augmented state vector. Applying the \textit{reciprocal convexity lemma} (Park et al., 2011), if there exists a matrix $S \in \mathbb{R}^{m \times m}$ such that $\begin{bmatrix} R & S \\ S^T & R \end{bmatrix} \ge 0$, the partitioned segments can be bounded. We obtain:
\begin{equation}
\dot{V}(t) \le -\eta^T(t) \Omega_{ISS} \eta(t) + \gamma \|d_t(t)\|^2
\end{equation}
where the ISS-gain parameter $\gamma$ is explicitly formulated as:
\begin{equation}
\gamma = \frac{1}{\epsilon_1}\lambda_{max}(P^2) + \tau_m^2 \left( 1 + \frac{1}{\epsilon_2} \right) \lambda_{max}(R)
\end{equation}
and the $3m \times 3m$ matrix $\Omega_{ISS}$ is defined as:
\begin{equation}
\Omega_{ISS} = \begin{bmatrix} 2 P K_p - Q - \epsilon_1 I_m & P K_d & 0 \\ * & (1-d)Q & 0 \\ * & * & 0 \end{bmatrix} + \begin{bmatrix} R & S-R & -S \\ * & 2R-S-S^T & S-R \\ * & * & R \end{bmatrix} - \tau_m^2 (1 + \epsilon_2) \begin{bmatrix} K_p^T R K_p & K_p^T R K_d & 0 \\ * & K_d^T R K_d & 0 \\ * & * & 0 \end{bmatrix}
\end{equation}

By solving the LMI $\Omega_{ISS} > 0$ under the constraint $\begin{bmatrix} R & S \\ S^T & R \end{bmatrix} \ge 0$, the derivative satisfies $\dot{V}(t) \le -\lambda_{min}(\Omega_{ISS}) \|\eta(t)\|^2 + \gamma \|d_t(t)\|^2$, proving the Local Input-to-State Stability (Local ISS) of the tracking error.

\subsection{Null-Space Projection for Joint Limit, Collision, and Posture Avoidance}
The secondary joint velocity vector $\dot{q}_0(t)$ is designed to push the robot joints away from physical boundaries, avoid self-collisions, and guide the redundant hand joints toward natural anthropomorphic configurations to resolve the $19 - 5 = 14$ DoFs under-constraint. We define the objective function $H(q) \in \mathbb{R}$ to be minimized:
\begin{equation}
H(q) = w_{lim} \sum_{i=1}^{n} \left( \frac{q_i - \bar{q}_i}{q_{i,max} - q_{i,min}} \right)^2 + H_{coll}(q) + w_{post} \sum_{j=7}^{25} \left( q_j - q_{j,nom} \right)^2
\end{equation}
where:
\begin{itemize}
\item $q_{i,max}, q_{i,min}$ are the physical boundaries of the $i$-th joint, and $\bar{q}_i = \frac{q_{i,max} + q_{i,min}}{2}$.
\item $q_{j,nom}$ denotes the nominal natural posture joint angle of the $j$-th hand joint, calibrated from the operator's relaxed hand position.
\item $w_{lim}$ and $w_{post}$ are positive weight parameters.
\end{itemize}

To prevent gradient explosion and chattering in discrete-time execution (100 Hz), the collision potential $H_{coll}(q)$ is active only within a finite influence distance $d_{influence} > 0$:
\begin{equation}
H_{coll}(q) = \sum_{a < b} \frac{\sigma}{2} \left( \frac{1}{d_{ab}(q)} - \frac{1}{d_{influence}} \right)^2 \quad \text{if } d_{ab}(q) \le d_{influence}, \text{ else } 0
\end{equation}
where $d_{ab}(q)$ is the distance between capsule representations of link $a$ and link $b$. 

To ensure the perturbation vector $d_l(t)$ remains bounded in the LKF analysis, the gradient mapping is saturated to bound the norm of the potential vector:
\begin{equation}
\dot{q}_0(t) = \text{sat}_{v_{max}} \left( -\alpha \nabla H(q) \right)
\end{equation}
where $\text{sat}_{v_{max}}(\cdot)$ is an element-wise saturation function bounding each component to $\pm v_{max}$. Consequently, the Euclidean norm of the secondary velocity is strictly bounded ($\| \dot{q}_0(t) \| \le \sqrt{n} v_{max}$). This guarantees that the DLS leakage perturbation $d_l(t)$ remains bounded globally, preserving the ISS conditions.

\textbf{Remark 2 (SVD-based Null-Space Leakage)}: Applying singular value decomposition (SVD) to the Jacobian $J = U \Sigma V^T$, the workspace leakage introduced by the DLS projection is given by $J \left( I_n - J^{\dagger}_{DLS}J \right) = U \text{diag}\left( \frac{\sigma_i \lambda^2}{\sigma_i^2 + \lambda^2} \right) V^T$. By AM-GM inequality, each diagonal entry is bounded by $\frac{\lambda}{2}$ for all $\sigma_i$. As the system approaches a deep singularity ($\sigma_i \to 0$), the leakage converges to zero. The maximum leakage is bounded globally by $\mathcal{O}(\lambda)$ at near-singular configurations ($\sigma_i \approx \lambda$), ensuring that the secondary tasks do not corrupt workspace accuracy.

\subsection{Practical Implementation Details}
In practice, the delay $\tau(t)$ is calculated in real-time using network packet timestamps: $\tau(t) = t_{recv} - t_{send}$. Because raw latency measurements contain high-frequency noise (jitter), the derivative $\dot{\tau}(t)$ cannot be computed directly via numerical differentiation. 

To resolve this issue, the controller implements a continuous-time Low-Pass Filter (LPF):
\begin{equation}
\tau_{filt}(s) = \frac{1}{T_f s + 1} \tau(s)
\end{equation}
To strictly guarantee Assumption 2 ($|\dot{\tau}_{filt}(t)| \le d$), the filter time constant $T_f$ cannot be chosen arbitrarily. Given the LPF time-domain dynamics $\dot{\tau}_{filt} = \frac{1}{T_f} (\tau - \tau_{filt})$ and the bounded jitter $\|n(t)\| \le n_{max}$ from Assumption 3, $T_f$ must be selected such that:
\begin{equation}
T_f \ge \frac{\max |\tau(t) - \tau_{filt}(t)|}{d}
\end{equation}
In practical implementation, $T_f$ is conservatively tuned above this theoretical lower bound to ensure smooth command velocities.

Furthermore, computing the collision gradient $\nabla H_{coll}(q)$ requires the real-time calculation of the closest-point Jacobians $J_{ab}(q)$ for all active capsule pairs:
\begin{equation}
\frac{\partial d_{ab}(q)}{\partial q} = \hat{n}_{ab}^T J_{ab}(q)
\end{equation}
where $\hat{n}_{ab}$ is the unit normal vector between the closest points. To ensure execution at 100 Hz, spatial hashing or bounding-volume hierarchy (BVH) algorithms provided by the Isaac Sim API are utilized to quickly cull inactive collision pairs ($d_{ab} > d_{influence}$), isolating the Jacobian computations only to the active proximity sets. This CLIK controller operates as a high-frequency (100 Hz) teleoperation bridge in the Isaac Sim simulator, allowing the human operator to generate stable, chatter-free demonstration data for downstream Diffusion Policy training.

\begin{thebibliography}{1}
\bibitem{ref1}
Y. Qin, W. Yang, B. Huang, and X. Wang, ``AnyTeleop: A General Vision-Based Dexterous Robot Arm-Hand Teleoperation System,'' in \textit{Proc. IEEE Int. Conf. Robot. Autom. (ICRA)}, 2023, pp. 10160--10166.
\bibitem{ref2}
J. Wang and L. Zhang, ``Event-Triggered Bilateral Teleoperation Control Under Time-Varying Communication Delays,'' \textit{IEEE Trans. Ind. Electron.}, vol. 69, no. 12, pp. 13245--13254, Dec. 2022.
\bibitem{ref3}
G. Palli, S. Salvatore, and C. Melchiorri, ``Bilateral Teleoperation of Asymmetric Robotic Systems with Time-Varying Delays,'' \textit{IEEE/ASME Trans. Mechatronics}, vol. 27, no. 4, pp. 2154--2165, Aug. 2022.
\bibitem{ref4}
X. Zhao, G. Zhang, and J. Wang, ``Event-triggered control and communication for single-master multi-slave teleoperation systems with Try-Once-Discard protocol,'' \textit{IEEE Trans. Autom. Control}, vol. 69, no. 3, pp. 1542--1549, Mar. 2024.
\bibitem{ref5}
H. B. Zhang, J. H. Park, and Y. Yoon, ``Novel Stability Criteria for Linear Time-Delay Systems Using Lyapunov-Krasovskii Functionals With A Cubic Polynomial on Time-Varying Delay,'' \textit{IEEE/CAA J. Autom. Sinica}, vol. 8, no. 5, pp. 1003--1011, May 2021.
\end{thebibliography}

\end{document}
